\documentclass[10pt,a4paper]{amsart}
\usepackage[margin=19mm]{geometry}
\usepackage{amsmath,amssymb,mathtools}
\usepackage[scr=rsfs,cal=euler]{mathalfa}
\usepackage{xfrac}
\usepackage[unicode,hidelinks,bookmarks=false]{hyperref}
\newcommand{\ie}{i.e.\ }
\newcommand{\cf}{cf.\ }
\newcommand{\resp}{respectively\ }
\newcommand{\Subsection}{\subsection}
\providecommand{\nobreakdash}{\nobreak\hbox{-}}
\setlength{\emergencystretch}{2em}
\multlinegap0pt
\usepackage{amssymb}
\usepackage{mathtools}
\usepackage{bm}
\usepackage{float}
\usepackage{xcolor}
\usepackage[shortlabels]{enumitem}
\usepackage{tikz}
\newtheorem{lemma}{Lemma}
\usepackage{cleveref}
\crefname{lemma}{Lemma}{Lemmas}
\newcommand{\E}{\mathbb{E}}
\newcommand{\delt}{\delta}
\newcommand{\eps}{\varepsilon}
\newcommand{\smallscale}{s_{\gamma,N}}
\title{Beyond Client Averaging: A Client-Independent Second-Order Stationary-Bias Component in Stochastic SCAFFOLD}
\author{Yi-Ping Tang, Guan-Ju Peng}
\date{Source: arXiv:2608.26765v1 (CC BY 4.0)}
\begin{document}
\maketitle

\noindent\emph{Source and licence.} Beyond Client Averaging: A Client-Independent Second-Order Stationary-Bias Component in Stochastic SCAFFOLD. Version arXiv:2608.26765v1 (CC BY 4.0), distributed by arXiv under the Creative Commons Attribution 4.0 International licence, http://creativecommons.org/licenses/by/4.0/, as recorded in the arXiv licence metadata for this exact version and shown on the abstract page https://arxiv.org/abs/2608.26765v1. Full licence notice: \texttt{licenses/arxiv-2608.26765-CC-BY-4.0.txt}.\par\smallskip

\noindent\textit{Editorial scope.} Counted unit: the article's lemma lem:Delta-moments (source0.tex lines 1294--1308). The article's complete original proof of it is delivered with it (source0.tex lines 1310--1359, one \texttt{proof} environment of 50 lines). No mathematics is written, repaired or re-derived: the statement and the proof are the article's own lines, in the article's own notation, and no retained line is substituted or re-flowed.  Retained uncounted as the source-local context the proof needs: lines 871--879; lines 976--983; lines 1212--1215; lines 1276--1281.  Nothing was omitted from any retained span, so the package carries no omission record.  No retained line contains a citation, so the wrapper carries no bibliography and no bibliography entry.  No retained line contains a figure environment, an image-inclusion command or drawing code, so no image is delivered and the package declares no asset dependency.  Editorial formatting only, in addition: an AMS \texttt{amsart} preamble that restates the article's own declarations and macros used by the retained lines, namely the environments \texttt{lemma} on separate counters, together with the standard packages those lines need.  The retained environments print in this document as Lemma 1 in order of appearance, whereas the article prints the same items with the numbering of its own full document.  The complete native source of the article as distributed in the arXiv e-print for this exact version is bundled unchanged as \texttt{sources/208658/source0.tex}.
\begin{lemma}[Uniform sixth control moment]
\label{lem:control-sixth}
There exist $\gamma_{0,H}^{(C6)}>0$ and $C_{\xi,6,H}<\infty$, independent of $N$ and $\gamma$, such that for $0<\gamma\le\gamma_{0,H}^{(C6)}$,
\begin{equation}
\label{eq:control-sixth}
  \frac1N\sum_{c=1}^N\E|\xi_c|^6\le C_{\xi,6,H}.
\end{equation}
Consequently, the client-averaged second and fourth control moments are uniformly bounded as well.
\end{lemma}

\begin{lemma}[Finite-step local displacement]
\label{lem:displacement2}
For every $h\in\{0,\ldots,H\}$,
\begin{equation}
\label{eq:disp2}
  \frac1N\sum_{c=1}^N\E[(\theta_{c,h}-x)^2]\le C_H\gamma^2.
\end{equation}
\end{lemma}

\begin{equation}
\label{eq:q-quadratic}
  |q(y)|\le \frac{K_3}{2}|y|^2.
\end{equation}

\begin{align}
  \eta&:=-\gamma\sum_{j=1}^H\lambda_\gamma^{H-j}\bar\eps_j,             \label{eq:eta-def}\\
  \Delta&:=-\gamma\sum_{h=0}^{H-1}\lambda_\gamma^{H-1-h}\bar q_h,
  \qquad
  \bar q_h:=\frac1N\sum_{c=1}^N q(\theta_{c,h}).                         \label{eq:Delta-def}
\end{align}

\begin{lemma}[Nonlinear-increment moments]
\label{lem:Delta-moments}
For the $\Delta$ in Equation~\eqref{eq:Delta-def},
\begin{align}
  \E[\Delta^2]&\le C_H\gamma^2\smallscale,                              \label{eq:Delta2}\\
  \E[\Delta^4]&\le C_H\gamma^7,                                         \label{eq:Delta4-sharp}\\
  \E|\Delta|^6&\le C_H\gamma^9.                                         \label{eq:Delta6}
\end{align}
A second, weaker but sometimes useful estimate is
\begin{equation}
\label{eq:Delta4-M4}
  \E[\Delta^4]\le C_H\gamma^4(M_4+\gamma^4),
  \qquad M_4:=\E[x^4].
\end{equation}
\end{lemma}

\begin{proof}
Besides the quadratic bound in Equation~\eqref{eq:q-quadratic}, smoothness gives the global linear bound
\[
  |q(y)|\le (L+a)|y|.
\]
The earlier coarse local fourth- and sixth-moment bounds give, uniformly for $h\le H$,
\[
  \frac1N\sum_c\E|\theta_{c,h}|^4\le C_H\gamma^2,
  \qquad
  \frac1N\sum_c\E|\theta_{c,h}|^6\le C_H\gamma^3.
\]
Using the quadratic bound, Jensen, and fixed $H$ yields $\E[\Delta^2]\le C_H\gamma^4$, which is stronger than Equation~\eqref{eq:Delta2} because $\smallscale\ge\gamma^2$.

For the fourth moment, combine the linear and quadratic bounds to obtain
\[
  |q(y)|^4\le C|y|^6.
\]
Jensen then yields
\[
  \E[\Delta^4]\le C_H\gamma^4
  \sum_{h=0}^{H-1}\frac1N\sum_c\E|\theta_{c,h}|^6
  \le C_H\gamma^7,
\]
which is the sharpened estimate needed for the fourth-moment absorption argument.

For completeness, we now prove the weaker estimate in Equation~\eqref{eq:Delta4-M4} without appealing to a later lemma. The same finite-step argument used in \cref{lem:displacement2} gives pathwise, for $\delt_{c,h}:=\theta_{c,h}-x$ and fixed $h\le H$,
\[
  |\delt_{c,h}|
  \le C_H\gamma\left(|x|+|\xi_c|+\sum_{j=1}^h|\eps_{c,j}|\right).
\]
Raise this inequality to the fourth power, average over clients and expectations, and use the earlier coarse global fourth-moment bound, the client-averaged fourth control moment from \cref{lem:control-sixth}, and bounded noise. This yields
\[
  \frac1N\sum_c\E|\delt_{c,h}|^4\le C_H\gamma^4.
\]
Consequently,
\[
  \frac1N\sum_c\E|\theta_{c,h}|^4
  \le 8M_4+8\frac1N\sum_c\E|\delt_{c,h}|^4
  \le C_H(M_4+\gamma^4).
\]
Using the global linear bound $|q(y)|\le (L+a)|y|$, Jensen over clients and the fixed sum over local steps now gives
\[
  \E[\Delta^4]
  \le C_H\gamma^4\sum_{h=0}^{H-1}\frac1N\sum_c\E|\theta_{c,h}|^4
  \le C_H\gamma^4(M_4+\gamma^4),
\]
which is exactly Equation~\eqref{eq:Delta4-M4}. Thus no forward reference to the later displacement-fourth-moment lemma is needed.

Finally, $|q(y)|^6\le C|y|^6$ by the linear bound, so the same argument gives Equation~\eqref{eq:Delta6}.
\end{proof}

\end{document}
